% from ega4-2-fr.tex
\oldpage[IV-2]{45}
$f(x)=g(y)=s$, soit $T_{x,y}$ le préschéma
$\operatorname{Spec}(k(x)\otimes_{k(s)}k(y))$, et soit $I_{x,y}$ l'image de
$\operatorname{Ass}(\mathcal O_{T_{x,y}})$ par le monomorphisme canonique
$T_{x,y}\to Z$ (I, 3.4.9). On a alors
\end{proposition}

\begin{env}[3.3.6.1]
\label{IV.2.3.3.6.1-fr}
\[
\operatorname{Ass}(\mathcal F\otimes_S\mathcal G)
=\bigcup_{x\in\operatorname{Ass}(\mathcal F)}
\left(\bigcup_{y\in\operatorname{Ass}(\mathcal G_{f(x)})}I_{x,y}\right),
\]
\end{env}

où pour tout $s\in S$, $\mathcal G_s=\mathcal G\otimes_{\mathcal O_S}k(s)$.

Soient $p:Z\to X$, $q:Z\to Y$ les projections canoniques, de sorte que l'on
a le diagramme commutatif
\[
\xymatrix{
X \ar[d]_f & Z \ar[l]_p \ar[d]^q \\
S & Y \ar[l]_g
}
\]
Posons $\mathcal G'=q^*(\mathcal G)$, de sorte que
$\mathcal F\otimes_S\mathcal G=\mathcal F\otimes_X\mathcal G'$~; comme
$\mathcal G'$ est $p$-plat (2.1.4), il résulte de (3.3.1) que l'on a

\begin{env}[3.3.6.2]
\label{IV.2.3.3.6.2-fr}
\[
\operatorname{Ass}(\mathcal F\otimes_S\mathcal G)
=\bigcup_{x\in\operatorname{Ass}(\mathcal F)}
\operatorname{Ass}(\mathcal G'_x)
\]
\end{env}

avec $\mathcal G'_x=\mathcal G'\otimes_{\mathcal O_X}k(x)$. Si $s=f(x)$, on a
$\mathcal G'_x=\mathcal G_s\otimes_{k(s)}k(x)$ et
$p^{-1}(x)=g^{-1}(s)\otimes_{k(s)}k(x)$~; en outre, comme le corps $k(x)$ est
un $k(s)$-module plat, le morphisme
$p^{-1}(x)\to g^{-1}(s)$ est plat (2.1.4)~; appliquant (3.3.1) à ce morphisme,
il vient

\begin{env}[3.3.6.3]
\label{IV.2.3.3.6.3-fr}
\[
\operatorname{Ass}(\mathcal G'_x)
=\bigcup_{y\in\operatorname{Ass}(\mathcal G_s)}
\operatorname{Ass}(\mathcal O_{T_{x,y}}),
\]
\end{env}

d'où la proposition.

On notera que si, dans l'énoncé, on supprime l'hypothèse
2\textsuperscript{o} on peut encore conclure, en vertu de (3.3.1), la relation

\begin{env}[3.3.6.4]
\label{IV.2.3.3.6.4-fr}
\[
\operatorname{Ass}(\mathcal F\otimes_S\mathcal G)
\supset\bigcup_{x\in\operatorname{Ass}(\mathcal F)}
\left(\bigcup_{y\in\operatorname{Ass}(\mathcal G_{f(x)})}I_{x,y}\right).
\]
\end{env}

\begin{corollary}[3.3.7]
\label{IV.2.3.3.7-fr}
Sous les hypothèses de (3.3.6), supposons en outre que $S$ soit localement
noethérien et que
$f(\operatorname{Ass}(\mathcal F))\subset\operatorname{Ass}(\mathcal O_S)$.
Alors on a
\end{corollary}

\begin{env}[3.3.7.1]
\label{IV.2.3.3.7.1-fr}
\[
\operatorname{Ass}(\mathcal F\otimes_S\mathcal G)
=\bigcup_{(x,y)\in C}I_{x,y},
\]
\end{env}

où $C$ est l'ensemble des couples $(x,y)$ tels que
$x\in\operatorname{Ass}(\mathcal F)$, $y\in\operatorname{Ass}(\mathcal G)$ et
$f(x)=g(y)$.

Comme $\mathcal G$ est $g$-plat, il résulte en effet de (3.3.1) que la relation
«~$s\in\operatorname{Ass}(\mathcal O_S)$ et
$y\in\operatorname{Ass}(\mathcal G_s)$~» équivaut à
$y\in\operatorname{Ass}(\mathcal G)$~; la conclusion résulte de (3.3.6.1).

\begin{remark}[3.3.8]
\label{IV.2.3.3.8-fr}
Nous verrons plus loin (4.2.2) que sous les hypothèses de (3.3.6),
$T_{x,y}$ est un préschéma \emph{sans cycle premier associé immergé}~; il en
résultera que si $\mathcal F$ et les $\mathcal G_s$ sont sans cycle premier
associé immergé, il en est de même de $\mathcal F\otimes_S\mathcal G$.
\end{remark}

\begin{corollary}[3.3.9]
\label{IV.2.3.3.9-fr}
Sous les conditions de (3.3.7), on a
\end{corollary}

\begin{env}[3.3.9.1]
\label{IV.2.3.3.9.1-fr}
\[
q(\operatorname{Ass}(\mathcal F\otimes_S\mathcal G))
\subset\operatorname{Ass}(\mathcal G)
\]
\end{env}

(où $q:X\times_S Y\to Y$ est la projection canonique).

En effet, si $(x,y)\in Z$, on a
$q(I_{x,y})=\{y\}\subset\operatorname{Ass}(\mathcal G)$.

